\section{Background}\label{sec:background}

\subsection{Unmanned Vehicles in Search and Rescue}\label{subsec:unmanned-vehicles-in-search-and-rescue}

Unmanned vehicles have long been used as force multipliers in SAR situations.
\cite{waharte2010supporting} shows that autonomous use of unmanned aerial vehicles (UAVs) could significantly enhance terrestrial SAR through structured search strategies, reducing expected detection times compared to random search.
In the maritime domain, recent SAR frameworks (\cite{messmer2024uavassistedmaritimesearchrescue}) for UAV-assisted MSAR have proposed software architectures that combine onboard UAV detection with more advanced ground-station processing.

\cite{7761074} demonstrated a co-operative multi-robot team for surveilling shipwreck survivors at sea, in collaboration with the Portuguese Navy, at an undeclared location.
While the paper offered promising results for the application of UAVs and unmanned surface vessels (USVs) in Maritime Search and Rescue, many such deployments retain strong human involvement at the control level.
In the experiment, a UAV takes off from the USV, ascends to a given height, and then pings a USV with the location of anything it suspects to be a person at sea.
The USV then moves to the target, and takes a picture of it, which is then sent with the target's location to a human operator for verification.

While this system is effective, it relies on a human operator, or a team of operators, to manage the drones in small groups.

This reliance on manual control helps to explain why many implementations of UAVs in MSAR do not use entirely autonomous solutions.
Factors like safety, certification, and legal accountability inherently favour conservative, operator-centric control approaches.
Consequently, there is little empirical evidence on how far autonomy and human-swarm interfaces can be pushed while maintaining acceptable workload and safety during a Maritime Search and Rescue operation.

\subsection{Swarm Coordination}\label{subsec:swarm-coordination}

Multi-agent swarm coordination is a well-studied topic which aims to achieve a robust method of coordinating multiple agents without single points of failure.
\cite{ENMASCA2025} introduced ``EN-MASCA'' (Enhanced Multi-Agent Swarm Control Algorithm) for drone swarm patrolling in durian orchards.
It employs reinforcement learning approaches to optimise formations and avoid obstacles in terrestrial agricultural settings.

While not related to the maritime domain, EN-MASCA shows that learned coordination can be deployed practically.
However, this comes with the drawback of requiring substantial training data, and computation resources.
Both of these issues can be prohibitive for SAR platforms.

Consensus-based approaches, like the one proposed by~\cite{dev2025swarnraft}, offer a more lightweight alternative to reinforcement learning.
A consensus-based approach tackles robustness to communication issues by allowing multiple drones to converge on a common estimate of the swarm state rather than rely on excessive ground-station communication.

Across these approaches, the most primary motivators appear to be collision avoidance and maximising search coverage.
Operators over open ocean may not need to worry about terrain, but care must be given to avoid obstacles like ships, waves, and other drones.

The Swarm Coordination method chosen was based on a set of criteria, which were informed by the project's goals of explainability, operator workload reduction, and practical constraints.
The criteria are as follows:

\begin{itemize}
    \item Operator Interpretability under time pressure
    \item Low computational overhead for near-real-time replanning
    \item Resilience to intermittent communications
    \item Balanced area coverage across heterogeneous drones
\end{itemize}

Several alternatives were considered.
EN-MASCA-style reinforcement learning-based approaches were discounted for this project phase since they require substantial training data, tuning effort, and additional validation to ensure its behaviour remains predictable for operators in safety-critical scenarios.
Pure consensus-based coordination algorithms like SwarmRaft (\cite{Dev_2026}) were also considered, but were discounted as the primary planner because, while robust to communication loss, it does not directly provide explicit spatial partitioning and workload balancing in a way that could be immediately visible in the interface.

Based on these constraints, BEACON adopted a geometric, explainable method.
The implemented solution combines weighted Voronoi tessellation to generate map ``cells'' assigned by drone position and speed (to encourage similar completion times), with Boustrophedon (lawnmower) waypoint generation inside each cell.
This design is less adaptive than learned policies in highly dynamic settings, but its behaviour is transparent and auditable, which better matches the human-supervisory and workload goals of BEACON\@.
The implementation of other coordination methods, while still possible, has been left for future work, due to the simulation-only scope of this project.

\subsection{Digital Twins}\label{subsec:digital-twins}
Digital Twins are commonly described as digital representations of physical assets or processes that remain synchronised with its real-world counterpart.
These are, notably, different from Digital Models (no data link), Digital Shadows (one-way physical-to-digital data link).

Digital Twins (DTs) can reduce cognitive workload by abstracting away low-level details, outsourcing tasks like aggregating different sensor streams/assorted metrics, enable predictive analysis, and supporting decision-making tools like path planners and alert systems.
These tools can then act consistently and transparently on behalf of the human supervisor.
In maritime applications, DTs have been used to optimise ship routes, monitor vessel health via operational metrics, with a focus on engineering and maintenance applications~\cite{DT_maritime_applications}.

BEACON adopts the ideas of DT in its layered architecture and decision-support functions, but it does not implement a real-time data link with physical drones, and instead simulates the drone swarm within the application itself.
This decision was made to focus on the design of the interface, and the swarm coordination algorithms, without needing to worry about the technical challenges of real-time data synchronisation, which will be implemented in future work.
It's therefore best described as a ``Digital Model/Shadow'' prototype, in that it's a high-fidelity simulation of a physical system with the accompanying User Interface (UI) that is designed to be part of a full Digital Twin system.
Being explicit about this scoping decision is important, to avoid conflating simulation with operational deployment and to focus on human-operator workload mitigation and the design of the human-swarm interface.

\subsection{HCI in Supervisory Control Systems}\label{subsec:supervisory-hci}

Human-computer interaction (HCI) research for supervisory control systems focuses on designing interfaces for human operators who oversee and manage autonomous systems.
\cite{AUTOMATIONEFFECTONUAVROUTING} found that excessively automated, simplistic displays can inadvertently increase operator workload by making debugging difficult, reducing situational awareness, and encouraging automation bias.
\cite{CUMMINGS2007339} found that operators supervising multiple UAVs, while experiencing high workload, would adopt load-shedding strategies, where they compartmentalise their attention to a small set of drones.

BEACON's interface combines an emphasis on revealing structure with constraints from research on supervisory control that warn against overly-complex global views of the system.
The mission map and coverage overlays act as configurable displays of the system's spatial constraints, like the mission boundary, Voronoi cells, and drone paths, while alerts and on-map notifications provide localised information.
The layout's design would be guided by various display design principles, like the Proximity Compatibility Principle.

\subsection{Workload Measurement and NASA-TLX}\label{subsec:workload-measurement-and-nasa-tlx}

A vital requirement for BEACON was a method to evaluate operator workload, to ensure that the system's design effectively reduces cognitive load while maintaining situational awareness.
The strategy for selecting a workload measurement was defined by ensuring suitability for repeated simulation trials, comparability with prior supervisory-control literature, relevance to the specific types of pressure common MSAR-like tasks (like cognitive, or temporal demand), and ease of administration.

Alternative strategies were considered, such as relying only on objective operational metrics (for example, intervention count and alert burden).
These strategies were not chosen as the primary method because objective metrics alone do not fully capture perceived strain.

Based on this selection process, NASA-TLX (Task Load Index) was adopted as the primary subjective workload measurement instrument.
It measures perceived workload across a set of dimensions:

\begin{itemize}
    \item Mental Demand - How mentally demanding was the task?
    \item Physical Demand - How physically demanding was the task?
    \item Temporal Demand - How hurried or rushed was the pace of the task?
    \item Performance - How successful were you in accomplishing what you were asked to do?
    \item Effort - How hard did you have to work to accomplish your level of performance?
    \item Frustration Level - How insecure, discouraged, irritated, stressed, and annoyed were you?
\end{itemize}

Each dimension is rated on a scale from 0 to 100, and then weighted based on pairwise comparisons to determine the relative importance of one dimension over another.

Recent reviews of NASA-TLX, however, have raised concerns about NASA-TLX's conceptual clarity, and its validity in interactive applications.
Systematic analyses of NASA-TLX have found that many studies use it without clearly defining workload, and apply inconsistent weightings, leading to difficulties in reproducing and comparing the results.

BEACON adopts NASA-TLX despite these caveats due to its widespread use and acceptance in human workload assessment, and its clean mapping to MSAR-like supervision, where mental and temporal demand, effort, and frustration are particularly important.
It is also practical for repeated use in simulation, and can be easily implemented.
At the same time, TLX scores are interpreted cautiously and are always complemented with objective operational metrics like alert burden, and intervention rate, in line with the recommendations given by~\cite{BABAEI2025103515}.
